% !TeX program = xelatex
% ============================================================================
%  第 6 课　课堂单页（学生用，单独打印一张 A4）
%  编译：XeLaTeX，两遍即可
%  记号：e 不动，r 转 120°，r² 转 240°，f 沿过 A 的轴翻
%  约定：复合 xy 读作「先做 x，再做 y」；格子里的动作是「先横排的，再做竖排的」
% ============================================================================

\documentclass[11pt]{ctexart}
\usepackage[a4paper,left=1.7cm,right=1.7cm,top=1.4cm,bottom=1.4cm]{geometry}
\usepackage{amsmath,amssymb}
\usepackage{booktabs}
\usepackage{array}
\usepackage[dvipsnames]{xcolor}

\setCJKmainfont{SimSun}[BoldFont=SimHei,ItalicFont=KaiTi]
\IfFontExistsTF{TeX Gyre Pagella}{\setmainfont{TeX Gyre Pagella}}{\setmainfont{Latin Modern Roman}}

\definecolor{pagegreen}{RGB}{34,139,34}
\renewcommand{\arraystretch}{1.25}
\setlength{\parindent}{0pt}
\pagestyle{empty}

\newcommand{\biaoti}[2]{%
  {\color{pagegreen}\rule{\textwidth}{1.2pt}}\\[6pt]
  {\heiti\Large 第 #1 课\quad #2}\hfill{\small\songti 课堂单页}\\[6pt]
  {\color{pagegreen}\rule{\textwidth}{0.5pt}}\\[4pt]
  {\small 姓名\underline{\hspace{3.2cm}}\qquad 班级\underline{\hspace{2.4cm}}\qquad 座号\underline{\hspace{1.6cm}}}
  \vspace{0.4em}
}

\begin{document}

\biaoti{6}{复合作为运算：\(S_3\) 是个群}

\section*{一、先把六个动作的两行式补上}

\begin{center}
{\small
\begin{tabular}{@{}clc@{}}
  \toprule
  名字 & 动作 & 两行式 \\
  \midrule
  \(e\) & 不动 & \(\begin{pmatrix}A&B&C\\A&B&C\end{pmatrix}\) \\[4pt]
  \(r\) & 顺时针转 \(120^\circ\) & \(\begin{pmatrix}A&B&C\\B&C&A\end{pmatrix}\) \\[4pt]
  \(r^2\) & 顺时针转 \(240^\circ\) & \(\begin{pmatrix}A&B&C\\\hspace{1.2em}&\hspace{1.2em}&\hspace{1.2em}\end{pmatrix}\) \\[4pt]
  \(f\) & 沿过 \(A\) 的轴翻 & \(\begin{pmatrix}A&B&C\\A&C&B\end{pmatrix}\) \\[4pt]
  \(rf\) & 沿过 \underline{\hspace{1.2em}} 的轴翻 & \(\begin{pmatrix}A&B&C\\C&B&A\end{pmatrix}\) \\[4pt]
  \(r^2f\) & 沿过 \underline{\hspace{1.2em}} 的轴翻 & \(\begin{pmatrix}A&B&C\\\hspace{1.2em}&\hspace{1.2em}&\hspace{1.2em}\end{pmatrix}\) \\
  \bottomrule
\end{tabular}}
\end{center}

\section*{二、乘法表：只填 \(r\) 行与 \(f\) 行}

\noindent 格子里的动作是「\textbf{先横排的，再做竖排的}」。

\begin{center}
{\small
\begin{tabular}{@{}l|cccccc@{}}
  \toprule
   & \(e\) & \(r\) & \(r^2\) & \(f\) & \(rf\) & \(r^2f\) \\
  \midrule
  \rule{0pt}{3.2ex}先 \(r\) & & & & & & \\
  \rule{0pt}{3.2ex}先 \(f\) & & & & & & \\
  \rule{0pt}{3.2ex}先 \(r^2\) & & & & & & \\
  \bottomrule
\end{tabular}}
\end{center}

\noindent 算不动的时候，就拿纸片真的做两遍，或者用两行式一格一格推。

\section*{三、四条规则，打勾}

\vspace{0.4em}
\noindent\begin{tabular}{@{}p{2.6cm}p{12.6cm}@{}}
  封闭 & \(\square\) 成立 \quad \(\square\) 不成立 \qquad 理由：\underline{\hspace{5.4cm}} \\
  \rule{0pt}{2.6ex}结合 & \(\square\) 成立 \quad \(\square\) 不成立 \qquad 理由：\underline{\hspace{5.4cm}} \\
  \rule{0pt}{2.6ex}单位元 & 它是 \underline{\hspace{2.2cm}} \\
  \rule{0pt}{2.6ex}逆元 & \(r\) 的逆是 \underline{\hspace{1.4cm}}，\(f\) 的逆是 \underline{\hspace{1.4cm}} \\
\end{tabular}

\vspace{0.6em}
\noindent 看一看 \(f\) 那一行的第二个格子，写出一条关系：\underline{\hspace{6cm}}

\end{document}